\section{Las dos comparaciones sobre historias detenidas}\label{historias}
Sea $\tau$ una parada que usa solamente el orden de filas y lo revelado hasta el paso actual.
Supongamos que cada padre anterior a la parada y cada uno de sus hijos tienen un predictor finito.
Escribimos $B_t$ para el tablero de la historia tras $t$ inserciones e $i_t$ para la fila que se
revela a continuación. Escribimos $\omega$ para el prefijo de casillas, incluyendo la última
inserción, $B_\tau$ para su residual y $\mu^\tau,\nu^\tau$ para las leyes conjuntas del orden y del
prefijo. Definimos
\[
 A_\tau=\sum_{t<\tau}\log\rho(B_t,i_t),\qquad g(B)=\log Q(B)-L(B)\quad(Q(B)>0).
\]

\begin{lema}[Comparación en la primera parada]\label{hist:igualdad}
Si $Q(B_0)>0$, se tiene $\mu^\tau\ll\nu^\tau$ y
\begin{equation}\label{super}
 g(B_0)=\E_\mu[A_\tau+g(B_\tau)]-D(\mu^\tau\Vert\nu^\tau).
\end{equation}
Aquí $D$ es la entropía relativa de las dos leyes discretas.
\end{lema}
\begin{proof}
Fijado el orden $\sigma$, cada una de las $Q(B_\tau)$ completaciones prolonga el prefijo de
una sola manera. Bajo $\nu$ multiplicamos las probabilidades de escoger las sucesivas casillas.
Las diferencias de los predictores se suman telescópicamente. Así, las masas en el mismo átomo son
\[
 \mu^\tau(\sigma,\omega)=\frac{Q(B_\tau)}{n!\,Q(B_0)},\qquad
 \nu^\tau(\sigma,\omega)=\frac{\exp\{L(B_\tau)-L(B_0)-A_\tau\}}{n!}.
\]
La segunda masa es positiva en el soporte de la primera. Tomar el logaritmo de su cociente
en el orden $\mu^\tau/\nu^\tau$ da $A_\tau+g(B_\tau)-g(B_0)$. Su media bajo $\mu^\tau$
es, por definición, la entropía relativa que figura en \eqref{super}.
\end{proof}

\begin{proposicion}[Cota inferior por permanencia]\label{hist:inferior}
Sea $\tau$ la parada de la Definición~\ref{def:regular}: la primera salida o, si no la hay, la
llegada al tamaño $K$. Sea $F$ el suceso de permanencia. En $F$ se tiene $B^{(K)}=B_\tau$.
Sea $G$ el conjunto de las historias de $F$ con $Q(B^{(K)})>0$. Si $s=\nu(G)>0$, entonces
\begin{equation}\label{infer}
 g(B_0)\ge \E_{\nu(\cdot\mid G)}[A_\tau+g(B^{(K)})]+\log s.
\end{equation}
\end{proposicion}
\begin{proof}
Para cada orden, los prefijos detenidos distintos tienen conjuntos disjuntos de completaciones.
En cada historia de $G$, multiplicar su masa bajo $\nu^\tau$ por $\exp\{A_\tau+g(B^{(K)})\}$
deja $Q(B^{(K)})e^{-L(B_0)}/n!$. Sumando primero en los prefijos y después en los órdenes obtenemos
\[
 \int_G\exp\{A_\tau+g(B^{(K)})\}\,d\nu^\tau\le Q(B_0)e^{-L(B_0)} .
\]
La integral es $s$ veces la esperanza condicional del exponencial. La concavidad del logaritmo
da \eqref{infer}, también cuando el condicionamiento modifica la distribución del orden.
\end{proof}

\begin{corolario}[Consistencia acumulada]\label{hist:consistencia}
En cada padre anterior a la parada supongamos las hipótesis del Teorema~7.1 y del Corolario~7.2
de~\cite{pred} para $k>M=\lceil\delta n\rceil$, y las del Teorema~8.1 del mismo trabajo para
$K<k\le M$. Entonces $|A_\tau|\le C/K+O_\delta(n^{-1})$ uniformemente en las historias.
\end{corolario}
\begin{proof}
Esos resultados dan $|\log\rho(B,i)|\le C_*/k^2$ en cada fila, incluyendo en $\rho$ todos los
hijos de esa fila. Para los padres con $k>M$ la constante depende de $\delta$ y de las cotas de
difusión, vacancia y brecha. Para los padres con $K<k\le M$ depende de las constantes de la clase
débil de~\cite[Sección~8.1]{pred}. En este rango, un padre regular de la Definición~\ref{def:regular}
está en la clase interior, por definición si $k<M$ y por la Proposición~\ref{meso:reconstruccion}
si $k=M$. La suma de $k^{-2}$ para $k>m$ es a lo sumo $\int_m^\infty t^{-2}\,dt=1/m$.
Una historia que termina antes usa un
subconjunto de estos sumandos.
\end{proof}
